\section{\foundry: Connections Instead of Declared Agents}
\label{sec:foundry}

The analysis reads agents out of a deployment, so it needs the deployment to
declare them. On \foundry it does not. We collected \nFoundryRepos
repositories, \nFoundryFiles files and \nFoundryDirs template directories, and
\nFoundryAgentDirs of those directories declare an agent as a deployed
resource. Agents there are created at run time through the service's own APIs,
so the deployment-time artifact the AWS analysis reads does not exist.

No rate comparable to the AWS one can be computed from these templates.

What \foundry templates do carry is \emph{connections}, and a connection's
\api{authType} decides whether a call presents a directory identity or a
credential stored in the connection. Where several agents share a project, a
stored credential is presented identically by all of them, so the target
cannot tell them apart and each agent reaches everything the connection is
trusted with. That is the same loss of separation the role sharing of
Section~\ref{sec:results} produces on AWS, reached by a different route.
Figure~\ref{fig:authtypes} gives the distribution:
\nAuthStoredCredential of \nFoundryAuthValues assigned values, about
\pctAuthStoredCredential percent, store a credential rather than presenting an
identity. \api{authType} is also a field on three unrelated Azure providers,
so the count is taken only inside directories declaring a connection resource.
No template sets an identity type scoped to an individual agent; those values
appear in parameter lists as allowed values only.

\begin{figure}[t]
\centering
\resizebox{\columnwidth}{!}{%
\begin{tikzpicture}
\begin{axis}[
  width=68mm,
  height=46mm,
  xbar,
  y dir=reverse,
  xmin=0, xmax=340,
  ymin=-0.7, ymax=6.7,
  bar width=3.4mm,
  /pgfplots/bar shift=0pt,
  axis x line*=bottom,
  axis y line*=left,
  y axis line style={draw=none},
  ytick={0,1,2,3,4,5,6},
  yticklabels={%
    \texttt{AAD},
    \texttt{ApiKey},
    \texttt{ManagedIdentity},
    \texttt{CustomKeys},
    \texttt{ProjectManagedIdentity},
    \texttt{AccessKey},
    seven others, one each},
  ytick style={draw=none},
  tick label style={font=\scriptsize},
  xtick={0,100,200,300},
  xlabel={\small assigned \texttt{authType} values},
  xlabel style={font=\small},
  nodes near coords={\scriptsize\pgfmathprintnumber\pgfplotspointmeta},
  every node near coord/.append style={anchor=west, xshift=0.5mm},
]
\addplot[xbar, bar shift=0pt, fill=cIdentity, draw=cIdentity] coordinates {(279,0)};
\addplot[xbar, bar shift=0pt, fill=cStored,   draw=cStored]   coordinates {(161,1)};
\addplot[xbar, bar shift=0pt, fill=cIdentity, draw=cIdentity] coordinates {(44,2)};
\addplot[xbar, bar shift=0pt, fill=cStored,   draw=cStored]   coordinates {(13,3)};
\addplot[xbar, bar shift=0pt, fill=cIdentity, draw=cIdentity] coordinates {(8,4)};
\addplot[xbar, bar shift=0pt, fill=cStored,   draw=cStored]   coordinates {(4,5)};
\addplot[xbar, bar shift=0pt, fill=cNeutral,  draw=cNeutral]  coordinates {(7,6)};
\end{axis}
\end{tikzpicture}}
\caption{\api{authType} across the \nFoundryAuthValues assigned values in the
\nFoundryAuthDirs directories that declare a connection resource. Blue
presents a directory identity, orange stores a credential in the connection.
The last row collects seven values that occur once each, which the artifact
names. \nAuthStoredCredential of \nFoundryAuthValues store a credential rather
than presenting an identity.}
\label{fig:authtypes}
\end{figure}

